\chapter{Theoretical Background}
\label{chap:fundamentacao}

The theoretical background is where you bring in the important concepts that underpin your work. You cite well-established authors, charts you found on Google Scholar at 2am, always citing the source with \enquote{According to the studies of \cite{Ding2020AdversarialAttacks}}.

This way, you build a text where your paragraphs are drawn directly from several works, showing that you aren't just making things up.

The secret here is to look like you dove deep into the literature, even if you really only dove into the Google Scholar "recommended" tab.

To do this, it helps to create several subtopics and sections. You can also include equations:

\begin{equation}
    x = \frac{-b \pm \sqrt{b^2 - 4ac}}{2a}
    \label{eq:quadratic}
\end{equation}

Or even code, where you have more freedom to bring in knowledge you elaborated yourself.

\lstinputlisting[language=Python,
caption=Example code
,label=lst:exemplocodigo]{content/\contentlangdir/chapters/fundamentacao/example.py}
\hfill
\begin{minipage}[t]{.65\textwidth}
\fonteautor
\end{minipage}
%example of inputs, ideal for organization and being able to reorder things later is to have one element per file
